% vim: tw=78 ai spell spelllang=cs

\documentclass{article}

\usepackage{graphicx}
\usepackage{hyperref}
\usepackage[utf8x]{inputenc}
\usepackage[czech]{babel}
\usepackage[T1]{fontenc}
\usepackage{xspace}

\pagestyle{empty}

\newcommand{\topinfo}[4]{
\begin{flushright}
\today
\end{flushright}
\vspace{1em}
\begin{center}
\Large\textbf{Posudek #1 na #2 práci \\
\normalsize%
\textit{#3:} #4}
\end{center}
}
\newcommand{\botinfo}{
\vspace{3em}
\begin{flushright}
Jiří Slabý
\end{flushright}
}


\newcommand{\cdb}{\textsc{ClabureDB}}

\begin{document}
\topinfo{vedoucího}{bakalářskou}{Viktor Hrtánek}{Podpora pro vkládání záznamů
do \cdb}

Dle zadání měla bakalářská práce nastudovat vkládání dat do databází různých
systémů a doimplementovat takovou funkcionalitu do fakultního nástroje \cdb,
a to s podporou dvou vstupních formátů, přičemž měla být korektně ošetřena
přístupová práva a celá implementace zdokumentována. Text práce je na 25
stranách strukturovaný do 6 kapitol.  V úvodu autor krátce popisuje motivaci
pro práci.  Následující kapitola popisuje různé databáze benchmarků. Další
kapitolou je kratinký popis systémů pro sledování chyb. Následuje kapitola s
popisem implementace a závěr. \textbf{Rozsah} práce se pohybuje mezi spodní
hranicí a doporučením požadavků na bakalářské práce.

\textbf{Obsahově} jsou dle mého názoru dostatečně propracovány všechny
kapitoly až na úvod. Ten by mohl lépe do práce uvádět a motivovat ji.  Ostatní
kapitoly by se ale stále daly rozšířit o užitečné informace např.\ z průběhu
implementace. Práce, stejně jako původní verze, obsahuje stále chyby.
Příkladem budiž následující:
\begin{itemize}
\item Příloha B: špatné příkazy. Všechny by měly být pomocí \texttt{bundle}.
V bodě 2 je navíc \texttt{exec} a navíc přípona \texttt{.rb}. Bez těchto změn
nemá běžný uživatel šanci systém zprovoznit.
\end{itemize}

Po \textbf{jazykové} stránce je práce již na průměrné úrovni. Bylo opraveno
spoustu chyb, ale objevily se nové. Např.\ ,,spolu s framework'', ,,zmenila
schéma databázy'' (chybí ,,sa''). V práci se objevuje stále několik
anglikanismů (,,Benchmark suites'').

\textbf{Praktická část} práce je v přiloženém \texttt{tar.gz} souboru. Protože
se jedná o rozšíření již existujícího nástroje \cdb, není v práci zcela přesně
uvedeno, co bylo změněno a přidáno. \textsc{Git} napovídá 24 souborů a 553
změněných řádků, přičemž některé jsou generované. Tento nově přidaný kód ale
není zcela funkční. Zejména mám tyto připomínky:
\begin{itemize}
\item Heslo se při změně nezmění, jak je psáno v příloze B (beze změny od
minulého pokusu).
\item Mohu importovat jakýkoliv typ souboru (beze změny od minulého pokusu).
Pak ho lze stáhnout a prohlížet (např.\ PDF). Nicméně při stisknutí
,,Approve'' pak aplikace samozřejmě přestane správně pracovat a objeví se šedá
obrazovka s výpisy.
\item ,,Approve'' jak textového, tak sqlite souboru selže s hlášením:
\begin{verbatim}
SQLite3::ConstraintException: NOT NULL constraint failed:
error.user: INSERTINTO "error" DEFAULT VALUES}
\end{verbatim}
\end{itemize}

Na závěr konstatuji, že přes všechny uvedené připomínky práce \textbf{splňuje
zadání a požadavky na bakalářskou práci}. Musím uvést, že autor pracoval sám a
o odevzdání dalšího pokusu jsem se dozvěděl až od vedoucího jedné z kateder
FI. Doporučuji práci \textbf{uznat} jako bakalářskou, ale vzhledem k výše
uvedeným okolnostem doporučuji hodnocení stupněm E.

\botinfo
\end{document}
